\section{Related Work}
\label{sec:related_work}

\paragraph{RL Training for LLM Agents.}
Group Relative Policy Optimization (GRPO)~\cite{shao2024deepseekmath} emerged from mathematical reasoning alignment and was scaled to large systems through techniques such as dynamic sampling and clip-higher stabilization~\cite{dapo2025}. Adapting GRPO to multi-turn interactive agents requires addressing the mismatch between its episode-level design and the practical constraints of long trajectories. RAGEN~\cite{ragen2025} introduced the step-independent windowing approach that has become the community default, partitioning trajectories into fixed-size windows optimized separately. Subsequent work has proposed algorithmic improvements that operate within or alongside this framework. GiGPO~\cite{feng2025gigpo} addresses credit assignment through a nested group structure that estimates step-level advantages without a learned critic, achieving substantial gains on ALFWorld and WebShop. Self-Guide~\cite{selfguide2026} co-evolves an internal reward model with the policy, providing denser learning signals than binary episode outcomes. SPEAR~\cite{spear2025} combines self-imitation learning with progressive curriculum to handle sparse rewards in early training. WebAgent-R1~\cite{wei2025webagent-r1} extends GRPO to multi-turn web interaction (M-GRPO) with dynamic context compression, while SORL~\cite{li2025sorl} introduces turn-level importance sampling for off-policy stability. HCAPO~\cite{hcapo2026} uses an LLM-as-critic to assign hindsight credit at critical decision states. These methods improve training algorithms or reward design at various granularities. None diagnoses the step-independent windowing mechanism itself as a source of systematic sign misattribution, which is the structural failure we identify and formalize.

\paragraph{Credit Assignment in Reinforcement Learning.}
The temporal credit assignment problem has a long history. REINFORCE~\cite{williams1992reinforce} provides unbiased but high-variance episode-level gradients. PPO~\cite{schulman2017ppo} reduces variance through a learned value function and clipped surrogate objectives. Generalized Advantage Estimation (GAE)~\cite{schulman2016gae} interpolates between one-step bootstrap and Monte Carlo returns via the $\lambda$ parameter, offering a continuous spectrum from TD(0) to full-episode credit. The TD($\lambda$) family~\cite{sutton1988td} established this interpolation in the tabular setting. In the LLM agent context, process reward models provide step-level supervision: AgentPRM~\cite{choudhury2025agentprm,xi2025agentprm} estimates per-step value through Monte Carlo rollouts or TD-based estimation with GAE. SkillMaster~\cite{skillmaster2026} proposes DualAdv-GRPO, which separately estimates advantages for task-solving and skill-editing decisions. Episode-level GRPO, as we advocate, is equivalent to GAE(1,1) without a critic: it assigns uniform advantage across all steps, forgoing per-step granularity in exchange for correct gradient direction. Our contribution is not a new credit assignment algorithm but a diagnosis of why the community's standard adaptation of GRPO (step-independent windowing) produces catastrophic sign errors in multi-turn settings.

\paragraph{Selective Verification and Adaptive Computation.}
A growing body of work trains models to allocate computation adaptively rather than uniformly. In the reasoning domain, Thinkless~\cite{du2025thinkless} introduces an RL-trained control token that decides whether to invoke extended thinking for each query. AutoThink~\cite{tu2025autothink} uses multi-stage RL to shape when reasoning-capable models activate their chain-of-thought. SelfBudgeter~\cite{selfbudgeter2026} frames adaptive computation as token budget allocation trained via budget-guided GRPO. In the agent domain, ARES~\cite{yang2026ares} trains an SFT-initialized router that selects reasoning effort levels per step, refined via GRPO with a combined outcome and token-efficiency reward. Agent-Omit~\cite{agent-omit2026} learns to omit redundant thoughts and observations through agentic RL, reducing context length without sacrificing performance. Sherlock~\cite{sherlock2025} performs selective verification at the workflow level through offline counterfactual analysis and speculative execution. Derailer-Rerailer~\cite{wan2025derailer-rerailer} detects instability triggers in reasoning chains and invokes targeted verification. These approaches operate at the reasoning-token level (Thinkless, AutoThink, SelfBudgeter), workflow-DAG level (Sherlock), or use heuristic triggers (Derailer-Rerailer) and SFT-initialized routers (ARES). Our selective verification differs in scope and mechanism: it makes a binary verify-or-skip decision at each agent interaction step, trained end-to-end via RL with an action-change informativeness signal that requires no external annotations or heuristic design.

\paragraph{Self-Verification and Self-Correction.}
Self-verification has received significant attention as a mechanism for improving LLM reliability. Kamoi et al.~\cite{kamoi2024self-correction-survey} survey self-correction in LLMs and establish that pure intrinsic self-correction (without external feedback) is unreliable across most settings. Long et al.~\cite{long2026selfverification} quantify the verification dilemma empirically: over 90\% of self-verification instances are confirmatory, with only approximately 2\% producing corrective changes. CorrectBench~\cite{tie2025correctbench} further demonstrates that uniform self-correction exhibits diminishing returns, as models repeatedly ``correct'' already-correct answers. VerifiAgent~\cite{verifiagent2025} proposes a unified verification framework with meta-verification and tool-based adaptive strategies, focusing on improving verification quality through better methods. These findings collectively establish two facts: verification quality is limited by confirmatory bias, and uniform application of verification wastes computation on uninformative instances. Our work addresses a complementary question: given that most verification is confirmatory, \textit{when} should an agent verify? The 90\% confirmatory rate reported by Long et al. directly motivates learning a selective policy that triggers verification only at steps where action change is likely.
